% from ega4-2-fr.tex
\oldpage[IV-2]{15}
  \item[(iii)] Pour toute partie fermée irréductible $Y'$ de $Y$, toute
  composante irréductible de $f^{-1}(Y')$ domine $Y'$.
\end{enumerate}
\end{proposition}

Il suffit par exemple de prouver (ii). Par hypothèse, il y a un
$\mathcal O_X$-Module quasi-cohérent de type fini $\mathscr F$, qui est
$f$-plat et tel que $\operatorname{Supp}(\mathscr F)=X$. Pour tout $x\in X$,
$\mathscr F_x$ est alors un $\mathcal O_x$-module de type fini, non réduit à
0, et en outre $\mathscr F_x$ est un $\mathcal O_y$-module \emph{plat}, pour
l'homomorphisme $\rho:\mathcal O_y\to\mathcal O_x$. Comme ce dernier est
local et que $\mathscr F_x\neq0$, le lemme de Nakayama montre que
$\mathscr F_x\otimes_{\mathcal O_y}k(y)\neq0$, donc $\mathscr F_x$ est un
$\mathcal O_y$-module \emph{fidèlement plat} $(0_{\mathrm I},6.4.1)$. Il
en résulte que pour tout idéal premier $\mathfrak q$ de $\mathcal O_y$, il
existe un idéal premier $\mathfrak p$ de $\mathcal O_x$ tel que
$\mathfrak q=\rho^{-1}(\mathfrak p)$ $(0_{\mathrm I},6.5.1)$, ce qui
prouve (ii).

\begin{corollary}[2.3.5]
\label{IV.2.3.5-fr}
Soit $f:X\to Y$ un morphisme vérifiant les conditions équivalentes (i), (ii),
(iii) de (2.3.4) (en particulier un morphisme quasi-plat ou un morphisme
ouvert (1.10.4)).
\begin{enumerate}
  \item[(i)] Soient $Z$, $Z'$ deux parties fermées irréductibles de $Y$ telles
  que $Z\subset Z'$, et soit $T$ une composante irréductible de $f^{-1}(Z)$~;
  alors il existe une composante irréductible $T'$ de $f^{-1}(Z')$ contenant
  $T$ (et dominant $Z'$).
  \item[(ii)] Pour toute composante irréductible $T$ de $X$,
  $\overline{f(T)}$ est une composante irréductible de $Y$.
  \item[(iii)] Supposons $Y$ irréductible, et, si $y$ est son point générique,
  supposons que $f^{-1}(y)$ soit irréductible. Alors $X$ est irréductible.
\end{enumerate}
\end{corollary}

(i) Il suffit d'appliquer (2.3.4, (i)) en prenant pour $x$ le point générique
de $T$ ($y=f(x)$ étant alors le point générique de $Z$) et pour $y'$ le point
générique de $Z'$.

(ii) Il est clair que $\overline{f(T)}=Z$ est irréductible, et en vertu de (i),
$Z$ ne peut être strictement contenue dans une partie fermée irréductible $Z'$
de $Y$.

(iii) En vertu de (ii), toute composante irréductible de $X$ domine $Y$, donc
rencontre $f^{-1}(y)$~; la conclusion résulte alors de
$(0_{\mathrm I},2.1.8)$.

\begin{proposition}[2.3.6]
\label{IV.2.3.6-fr}
Soit $Y$ un préschéma dont l'ensemble des composantes irréductibles est
\emph{localement fini} (ce qui est le cas si l'espace sous-jacent à $Y$ est
\emph{localement noethérien} (cf. I, 6.1.9)).
\begin{enumerate}
  \item[(i)] Toute partie fermée $W$ de $Y$ \emph{stable par générisation}
  $(0_{\mathrm I},2.1.2)$ est ouverte. En particulier, toute composante
  connexe de $Y$ est ouverte.
  \item[(ii)] Soit $f:X\to Y$ un morphisme fermé vérifiant en outre les
  conditions équivalentes (i), (ii), (iii) de (2.3.4) (ce qui sera le cas si
  $f$ est quasi-plat ou ouvert). Alors l'image par $f$ de toute composante
  connexe $C$ de $X$ est une composante connexe de $Y$.
  \item[(iii)] Soit $f:X\to Y$ un morphisme vérifiant les conditions
  équivalentes (i), (ii), (iii) de (2.3.4) et tel en outre que pour tout
  $y\in Y$, l'ensemble $f^{-1}(y)$ soit fini (ce qui sera le cas si $f$ est
  quasi-fini ou radiciel). Alors l'ensemble des composantes irréductibles de
  $X$ est localement fini.
\end{enumerate}
\end{proposition}

